% !TEX root = ../tree_notes.tex
% "How to Read This Book". This chapter is also the visual style guide:
% every box below is produced by the real definitions in setup/, so what you
% see here is exactly what chapters look like. Front matter uses starred
% (unnumbered) boxes and never calls \solvedproblem / \practice.
\chapter{How to Read This Book}

This book teaches the tree data structures and tree algorithms behind
competitive programming and coding interviews, in the order you need them.
Every idea is \emph{explained with a picture}, \emph{stated precisely},
\emph{turned into code}, and then \emph{used on real problems}. The same
layout repeats in all 31 chapters, so once you know how one topic is
organised you can navigate the whole book.

\section*{The shape of a chapter}

A chapter covers one area (for example \emph{Tree Traversals}) and is split
into \textbf{topics}. A topic is a numbered section such as~1.3, and it
always walks through the same nine stations, from memory-refresh to practice.

\begin{figure}[H]
\centering
\begin{tikzpicture}[diagram, node distance=4mm]
  \node[boxnode, fill=theorybg, draw=theoryborder] (op) {Chapter\\opener};
  \node[boxnode, right=6mm of op, fill=primaryblue!12, draw=primaryblue] (t1) {Topic\\N.1};
  \node[boxnode, right=of t1, fill=primaryblue!12, draw=primaryblue] (t2) {Topic\\N.2};
  \node[boxnode, right=of t2, fill=primaryblue!12, draw=primaryblue] (t3) {Topic\\N.3};
  \node[boxnode, right=of t3, fill=primaryblue!12, draw=primaryblue] (t4) {Topic\\N.4};
  \node[boxnode, right=6mm of t4, fill=successgreen!15, draw=successgreen] (end) {Chapter\\review};
  \draw[flow] (op) -- (t1); \draw[flow] (t1) -- (t2); \draw[flow] (t2) -- (t3);
  \draw[flow] (t3) -- (t4); \draw[flow] (t4) -- (end);

  \begin{scope}[yshift=-2.3cm]
    \foreach \i/\txt/\fl [count=\k] in {a/Previously\\(recall\\)/lightgray!60, b/Pattern\\card/white, c/Theory\\and code/white,
        d/Worked\\examples/white, e/Pitfalls\\WRONG vs RIGHT/white, f/Problem Lab\\(solved + ladder\\)/accentorange!18,
        g/Hardening\\AC $\to$ red/accentorange!18, h/Variations\\and twists/white, i/Practice\\problems/white} {
      \node[boxnode, minimum width=15mm, font=\footnotesize, fill=\fl] (s\i) at (\k*1.66-1.15,0) {\txt};
    }
    \foreach \a/\b in {a/b,b/c,c/d,d/e,e/f,f/g,g/h,h/i} \draw[softarrow] (s\a) -- (s\b);
  \end{scope}
  \draw[guide] (t2.south west) -- (sa.north west);
  \draw[guide] (t2.south east) -- (si.north east);
  \node[annot, below=3mm of sg.south east] {every topic follows the same nine stations};
\end{tikzpicture}
\caption*{A chapter is an opener, its topics and a 15-part review. Each topic
runs from a spaced-repetition refresh to graded practice in a fixed order.}
\end{figure}

\begin{itemize}
  \item The \textbf{opener} states what the chapter is for with a level badge on
        every promise, offers a \emph{beginner path} and a \emph{full path} through
        the material, draws the chapter on the six-step \emph{difficulty ladder},
        maps how the topics depend on each other, and ends with a
        \emph{problem recognition guide} ("when you see this, think that") and
        the notation table.
  \item Each \textbf{topic} opens with a \emph{Previously in this chapter} refresh
        and an \emph{In This Section} box, hands you a \emph{pattern card} for
        recognising the technique, builds the ideas with pictures and code, shows
        the classic trap as a WRONG-vs-RIGHT code pair, then runs a \textbf{Problem
        Lab}: an \emph{attempt gate} first, then fully solved problems walked step
        by step with a \emph{think it through} box (dead ends included) and an
        \emph{Approach Ladder} from the naive idea to the accepted one, finishing
        with adversarial \emph{hardening} notes, variations and a graded practice
        list. A \emph{recall card} with a self-test closes the section.
  \item The \textbf{chapter review} is fifteen fixed parts: summary tables, a
        master decision flowchart, a mistakes checklist, a \emph{WA clinic}
        (judge verdict $\to$ cause $\to$ where it was trained), a quick-reference
        card, write-from-memory drills, a complete-problem checklist, a study
        order, difficulty tiers, a reduction gallery, a mixed problem set, the
        chapter's summit, a cumulative tool-kit inventory, related chapters, and
        a last word.
\end{itemize}

\section*{The difficulty ladder}

Every promise, problem, warning and ladder rung in the book carries a level
badge $\lv1$--$\lv6$: $\lv1$ is \emph{what is a tree} (rating about 800),
$\lv3$ is the solid interview medium (about 1600), $\lv6$ is Codeforces
$2600$--$3500{+}$. A problem chip shows the official judge rating \emph{and}
the book level, e.g.\ \difficultytag{LC Medium L3}.

\begin{figure}[H]
\centering
\begin{tikzpicture}[diagram, x=9mm, y=7mm]
  \foreach \st/\lab in {1/lv1,2/lv2,3/lv3,4/lv4,5/lv5,6/lv6} {
    \fill[\lab!22, draw=\lab, line width=0.8pt] (\st-1,0) rectangle (\st,\st*0.52);
    \node[font=\scriptsize\bfseries, text=\lab] at (\st-0.5,\st*0.52-0.42) {L\st};}
  \node[annot, anchor=south west, font=\scriptsize] at (0.05,0.62) {foundations (Ch.\,1--2)};
  \node[annot, anchor=south west, font=\scriptsize] at (2.05,1.66) {interview core (Ch.\,3--12)};
  \node[annot, anchor=south west, font=\scriptsize] at (4.05,2.70) {contest depth (Ch.\,13--19)};
  \node[annot, anchor=south west, font=\scriptsize] at (5.05,3.22) {the summits (Ch.\,20--23)};
\end{tikzpicture}
\caption*{Six levels, one ladder, whole book. Each chapter's opener places its
topics on these steps and offers a beginner path ($\lv1$--$\lv3$) that skips
everything marked above it on a first pass.}
\end{figure}

\section*{The boxes}

Each coloured box has exactly one job, so you can skim a page by colour.

\begin{definition*}{Tree}{}
A \textbf{\textcolor{deepblue}{tree}} is a connected graph with no cycles. Equivalently, it is a set of nodes joined by edges where any two nodes are linked by exactly one path.\par
\emph{Blue boxes define a term precisely. The term is printed in bold blue where it is introduced.}
\end{definition*}

\begin{theorem*}{Edges of a tree}{}
A tree with $n$ nodes has exactly $n-1$ edges.\par
\emph{Purple boxes state results: theorems, lemmas (stepping stones) and corollaries (immediate consequences).}
\end{theorem*}

\begin{proofbox}
Every result is proved or argued, usually right below it. Proofs end with a black square. When a proof uses an earlier result it cites it by number, and the number is a link.
\end{proofbox}

\begin{insightbox}
Yellow boxes give the idea in one or two sentences: the picture to keep in your head.
\end{insightbox}

\begin{tracebox}
Green boxes walk through an algorithm by hand on a small tree, step by step, so you can check your understanding (and your code) against it.
\end{tracebox}

\begin{warningbox}
Red boxes flag the mistakes that cost accepted verdicts: forgetting the parent pointer and looping forever, off-by-one on depth vs.\ height, wrong \nil{} base case in recursion.
\end{warningbox}

\begin{complexitybox}
Grey boxes state running time and memory, and why.
\end{complexitybox}

\begin{notebox}[Note / Contest tip]
Light-grey boxes hold side remarks: history, alternative methods, judge-specific advice.
\end{notebox}

\begin{formulabox}
$\text{a full binary tree has } L = I + 1 \text{ leaves for } I \text{ internal nodes}$ \qquad a fact worth memorising
\end{formulabox}

\section*{The learning-loop boxes}

The new-structure boxes put the \emph{practice} side of reading into the book:
they make you retrieve, attempt, fail safely and harden -- shown here with the
real macros a chapter uses.

\begin{previouslybox}
\emph{Grey box at the top of every topic:} three lines that force you to recall
what the last topic built before a new idea is allowed to land on it.
\end{previouslybox}

\begin{recognitionbox}[Pattern card -- what triggers this technique]
\reckey{the exact words in a problem statement that mean ``use this'' -- read
them before any code.}
\recuse{
  \item the situation where the pattern is the fastest correct answer;
}\recnotuse{
  \item the look-alike that belongs to a later chapter -- named, with its pointer.
}
\end{recognitionbox}

\begin{attemptgate}
\emph{Orange gate before every Problem Lab:} the problems, their levels, and
where attempting them unaided is realistic. Only then do the solutions open.
\peek{One-line orientation you are allowed to peek at without spoiling it.}
\nudge{A stronger hint that names the tool but not the trick.}
\fullhint{The whole idea -- after reading this, code it yourself anyway.}
\end{attemptgate}

\begin{thinkbox}
\deadend{the obvious first idea}{the counterexample or complexity class that
kills it -- named, so you can reuse the autopsy.}
\aha{the one-sentence insight the whole solution hangs on.}
\end{thinkbox}

\begin{hardeningbox}
\emph{Purple box after an accepted solution:} an adversarial stress case, a
variant that breaks the assumptions, and a stretch that points at the next
chapter. $\lv3$ stress: try a $10^5$-node chain; \ldots
\end{hardeningbox}

\begin{recallbox}
\selftest{
  \item The blank you must fill \underline{\hspace{3cm}} from memory;
  \item one question per topic, ticked only when answered cold.
}
\emph{Closes every topic} -- tomorrow's revision starts here, not at page one.
\end{recallbox}

\section*{Extra tools for the tricky parts}

Some ideas -- rotations, binary lifting, heavy--light decomposition, lazy propagation --
are hard to grasp from prose alone. Four optional aids exist for exactly those moments.
They appear where an author needs them, not in every topic.

\begin{analogybox}[File system]
A rooted tree is a set of folders: the root is \code{/}, each folder is a node, its
sub-folders are children, and a folder with none is a leaf. ``The subtree at $v$'' is
everything inside one folder.\par
\emph{A blue analogy box ties an abstract definition to something you already know.}
\end{analogybox}

\usedin{\chref{c13} (segment trees) and \chref{c14} (Fenwick trees)}{the Euler tour that
flattens a subtree into a contiguous range is what turns ``sum over the subtree of $v$''
into a plain range query.}

\vsbox{Depth vs Height}{Depth $\depth(v)$}{Height $\height(v)$}{edges measured down
from the \emph{root} to $v$}{edges measured down from $v$ to its deepest \emph{leaf}}

\begin{figure}[H]
\centering
\begin{tikzpicture}[diagram]
  \node[tnode] (a1) at (0,0) {3};
  \node[tnode] (b1) at (-0.9,-1.3) {2};
  \draw[tedge] (a1) -- (b1);
  \node[tnode] (a2) at (3.4,0) {2};
  \node[tnode] (b2) at (4.3,-1.3) {3};
  \draw[tedge] (a2) -- (b2);
  \begin{scope}[on background layer]
    \node[frame,   fit=(a1)(b1)] (f1) {};
    \node[framehl, fit=(a2)(b2)] (f2) {};
  \end{scope}
  \node[stepno] at (f1.north west) {1};
  \node[stepno] at (f2.north west) {2};
  \draw[snap] (f1.east) -- (f2.west) node[midway, above, annot] {right rotate};
  \node[framecap, below=1mm of f1.south] {unbalanced};
  \node[framecap, below=1mm of f2.south] {balanced};
\end{tikzpicture}
\caption*{A \textbf{state sequence} shows an operation happening: numbered frames, a
thick orange arrow, and the key result boxed in orange. Harder practice items may also
carry a \textbf{hint ladder} -- a gentle nudge then a stronger one -- so you are never
stuck on a single clue.}
\end{figure}

\section*{Solved problems and practice problems}

A \textbf{solved problem} opens with a header like the one below and always
runs the same pipeline: \emph{Understanding the problem}, \emph{Key
observations}, \emph{Think it through} (two dead ends and the aha),
\emph{Solution}, \emph{Proof of correctness}, \emph{Implementation},
\emph{Complexity}, \emph{Approach Ladder} (every idea you could have had,
ranked, with the verdict each one earns), and \emph{Pitfalls and variants}
followed by the hardening box. The \texttt{WRONG}/\texttt{RIGHT} listings in
front of each lab show the chapter's classic bug as executable code.
Try the problem yourself before you read past the attempt gate.

\solvedheader{W1.1.1}{Maximum Depth of a Binary Tree}{LeetCode 104}{https://leetcode.com/problems/maximum-depth-of-binary-tree/}{LC Easy L1}{recursion, height of a tree}

\medskip
An \textbf{Approach Ladder} reads like this (the full one lives with each solved
problem): the level is what the idea costs you to invent, the verdict is what
the judge says.

\begin{center}
\renewcommand{\arraystretch}{1.3}
\begin{tabular}{@{}c p{0.235\linewidth} p{0.28\linewidth} l c@{}}
\toprule
\textbf{Level} & \textbf{Approach} & \textbf{Idea} & \textbf{Time} & \textbf{Verdict} \\
\midrule
$\lv1$ & BFS, count levels & queue the whole frontier & $O(n)$ & \ac \\
$\lv1$ & Recursive max-depth & one line, post-order & $O(n)$ & \ac \\
$\lv2$ & Height in edges at each node & off-by-one on the base case & $O(n)$ & {\colorbox{warningred!12}{\textcolor{warningred}{\scriptsize\bfseries WA}}} \\
\bottomrule
\end{tabular}
\end{center}

\medskip
\textbf{Practice problems} come right after each topic's lab, ordered from easy
to hard, each with a level badge and a one-line hint that nudges without
spoiling; the harder ones carry a second, stronger hint:

\begin{practiceset}
  \practiceitem{Q2.1.6}{Balanced Binary Tree}{LeetCode 110}{https://leetcode.com/problems/balanced-binary-tree/}{LC Easy L2}{Compute height bottom-up and return a sentinel (like $-1$) as soon as a subtree is unbalanced.}
\end{practiceset}

\section*{Numbering and cross-references}

\begin{center}
\begin{tabular}{@{}lll@{}}
\toprule
\textbf{Item} & \textbf{Looks like} & \textbf{Rule} \\
\midrule
Topic (section)                 & 1.3            & chapter.topic \\
Definition, theorem, \dots      & Theorem 1.3.4  & one shared counter per topic \\
Figure / table                  & Figure 1.3.2   & numbered per topic \\
Code listing                    & Code 1.3.1     & numbered per topic \\
Solved problem                  & W1.3.2         & \emph{W}orked, numbered per topic \\
Practice problem                & Q1.3.5         & \emph{Q}uestion, numbered per topic \\
Mixed problem (chapter review)  & M1.4           & numbered per chapter \\
\bottomrule
\end{tabular}
\end{center}

Every number is a clickable link in the PDF. The \emph{List of Solved
Problems} at the end of the book gives every W-problem with its source.

\section*{Difficulty tags}

Problems carry a tag. LeetCode uses its own labels; Codeforces ratings use
the Codeforces colour bands; other judges are shown in grey.

\begin{center}
\difficultytag{LC Easy}\ \difficultytag{LC Medium}\ \difficultytag{LC Hard}\qquad
\difficultytag{CF 900}\ \difficultytag{CF 1300}\ \difficultytag{CF 1500}\
\difficultytag{CF 1700}\ \difficultytag{CF 2000}\ \difficultytag{CF 2200}\
\difficultytag{CF 2600}\qquad
\difficultytag{CSES}\ \difficultytag{AtCoder 1580}
\end{center}

\section*{Notation}

\begin{center}
\begin{tabular}{@{}ll@{\qquad}ll@{}}
\toprule
$n$ & number of nodes & $\depth(v)$ & edges from the root to $v$ \\
$\height(v)$ & longest path down to a leaf & $\height(T)$ & $\height(\text{root})$ \\
$\size(v)$ & nodes in the subtree of $v$ & $\level(v)$ & $\depth(v)+1$ (1-based) \\
$\degop(v)$ & number of children / edges & $\parent(v)$ & the node above $v$ \\
$\children(v)$ & the nodes just below $v$ & $\lca(u,v)$ & lowest common ancestor \\
$\dist(u,v)$ & edges on the $u$--$v$ path & $\diam(T)$ & longest path in the tree \\
\bottomrule
\end{tabular}
\end{center}

We count depth and height in \emph{edges} unless a problem says otherwise;
when we count in nodes we say so. Inside $O(\cdot)$ the base of a logarithm
does not matter, so we write $O(\log n)$.

\section*{Code conventions}

\begin{itemize}
  \item All code is \cpp{} and is accepted by GCC with \code{-std=c++17} on
        Codeforces, AtCoder, CSES and the LeetCode compiler. Listings are
        complete: they compile as printed.
  \item Binary trees use the classic node struct with \code{left} and
        \code{right} pointers; \nil{} marks a missing child.
  \item General trees are stored as an adjacency list
        \code{vector<vector<int>>} and walked with DFS, always passing the
        parent so we never revisit it.
  \item Recursion depth can reach $n$; for a chain of $10^5$ nodes prefer an
        iterative traversal or raise the stack limit if the judge allows it.
\end{itemize}

\section*{How to practise with this book}

\begin{enumerate}
  \item Read the chapter opener's \emph{recognition guide} first; it decides
        which topic a future problem belongs to.
  \item In each topic, answer the \emph{Previously} and \emph{recall} cards from
        memory before reading on -- retrieval is the revision.
  \item At the \emph{attempt gate}, actually attempt: twenty minutes on the
        flagged problem before the lab opens, then one observation at a time.
        Use the peek, nudge and full-hint lines in that order, never as the
        first move.
  \item After an accepted solution, walk its \emph{Approach Ladder} aloud --
        say why each lower rung dies -- then read the \emph{hardening box} and
        run its stress case.
  \item Solve the practice list in order. Use the hint only after a real attempt.
  \item Close each chapter with the review's \emph{decision flowchart}, the
        \emph{write-from-memory drills} and the \emph{complete-problem
        checklist}; do the mixed set last -- those problems do not tell you
        which topic they belong to.
\end{enumerate}

\section*{How the notes are organised on disk}

The book is split so that each part can be written and compiled on its own:
\code{setup/} holds the whole preamble (packages, boxes, notation);
\code{chapters/chNN\_name/} holds one chapter, with one file per topic;
\code{structure/} holds the writing templates used for new chapters; and
\code{tools/preview.py} compiles a single topic in seconds while
\code{tools/build.py} builds the complete book. Problem counts in this book
(title page, part pages, chapter headers) are computed automatically.
